%% ============================================================
%%  Chapter 10 --- Qdrant Production Setup & Vector Store Architecture
%%  Production-Grade Agentic AI: LangGraph + Python Frameworks
%%  Dependencies: qdrant-client >= 1.11.0, mem0ai >= 0.1.0 (OSS)
%% ============================================================

\chapterquote{A vector store that ignores its own metadata is just an expensive random-number generator.}{anon}

Agent memory failures are silent. When an agent retrieves stale, irrelevant,
or cross-tenant memories, it does not crash --- it simply reasons from bad
evidence. A team running a multi-tenant assistant at fifty thousand users
discovered this the hard way when a misconfigured vector store began leaking
memories across user accounts. The root cause was not a bug in retrieval
logic: it was a missing payload index on the \code{user\_id} field, a field
that existed in every stored point but had never been registered with Qdrant.
Queries were falling back to full-collection scans, the HNSW graph was
navigating without tenant awareness, and the planner had no cardinality
estimate to guide its strategy. Recall had silently collapsed to near-zero
for most users while appearing functional in single-user integration tests.
This chapter prevents that failure mode.

The solution has four layers. The first layer is \emph{index geometry}: the
HNSW graph parameters that govern how the vector index is built and queried.
The second layer is \emph{metric commitment}: the distance function baked
permanently into the collection at creation time. The third layer is
\emph{payload indexing}: the structured metadata indexes that allow Qdrant's
query planner to make cardinality-aware decisions and enable filter-aware HNSW
traversal. The fourth layer is \emph{tenant architecture}: the collection
design pattern that keeps user data isolated without multiplying collection
overhead. These layers are not independent --- each depends on the previous one
being in place.

The sections in this chapter must be read in dependency order, and that order
is also the mandatory operational order:

\begin{enumerate}
  \item \textbf{Create the collection} with explicit HNSW config and distance metric (\S\ref{sec:hnsw-tuning} and \S\ref{sec:metric-commitment})
  \item \textbf{Create all payload indexes} before any data arrives (\S\ref{sec:payload-indexes})
  \item \textbf{Ingest data} --- only after steps 1 and 2 are complete
  \item \textbf{Query with filters} --- filter-aware HNSW traversal is only possible when indexes existed at segment build time (\S\ref{sec:filter-aware-hnsw})
\end{enumerate}

\begin{warningbox}
Reversing steps 2 and 3 does not raise an error. Qdrant will accept both data
and queries without complaint. The consequence is silent recall degradation
that is invisible in single-user tests and catastrophic at multi-tenant scale.
\end{warningbox}

Section~\ref{sec:multi-tenancy} applies the complete stack to the multi-tenant
isolation problem. Section~\ref{sec:mem0-backend} connects everything to the
mem0 OSS \code{Memory} class introduced in Chapter~8, closing the loop from
\S8.2.

\begin{table}[h]
\centering
\caption{Domain terms introduced in this chapter}
\label{tab:ch10-terms}
\begin{tabularx}{\textwidth}{lX}
\toprule
\textbf{Term} & \textbf{Definition} \\
\midrule
HNSW & Hierarchical Navigable Small World graph; the approximate nearest-neighbour index structure Qdrant builds over stored vectors. \\
\emph{ef} / \emph{ef\_construct} & Beam width during graph traversal: \code{ef\_construct} applies at index-build time, \code{ef} applies at query time. \\
Payload index & A field-level index on point metadata enabling fast filtering and cardinality estimation by the query planner. \\
Tenant index & A payload index variant (\code{is\_tenant=True}) that instructs Qdrant to physically co-locate data belonging to each tenant on disk. \\
Principal index & A payload index variant (\code{is\_principal=True}) that optimises range-based traversal on a field used as the primary ordering axis (e.g., a timestamp). \\
Filter cardinality & The estimated number of points that satisfy a filter predicate; used by the query planner to choose between filtered HNSW traversal and payload scan followed by reranking. \\
\bottomrule
\end{tabularx}
\end{table}

All code in this chapter lives under \code{src/vector\_store/} and
\code{src/memory/}. The shared \code{QdrantClient} instance is constructed
once in \code{src/vector\_store/client.py} and passed as an argument
everywhere else --- the client is never instantiated inside a function that
performs a search or upsert.

\begin{lstlisting}[language=bash, caption={Chapter source layout}, label={lst:ch10-layout}]
src/
  vector_store/
    client.py           # QdrantClient factory
    collection_setup.py # create_collection, bulk-load helpers
    index_setup.py      # create_payload_index loop
    memory_search.py    # search() and scroll() wrappers
    models.py           # MemoryPayload, MemoryPoint Pydantic models
    tenant_ops.py       # build_user_filter, upsert_memory, search_user_memories
  memory/
    mem0_client.py      # QdrantSettings, ensure_collection_ready, get_memory
\end{lstlisting}

The shared client factory is defined once and referenced throughout:

\begin{lstlisting}[language=Python,
  caption={Shared QdrantClient factory},
  label={lst:qdrant-client}]
# src/vector_store/client.py
from qdrant_client import QdrantClient
from src.memory.mem0_client import QdrantSettings

def get_qdrant_client(settings: QdrantSettings) -> QdrantClient:
    return QdrantClient(
        url=str(settings.url),
        api_key=settings.api_key.get_secret_value()
               if settings.api_key else None,
    )
\end{lstlisting}


%% ------------------------------------------------------------
\section{Index Geometry Is a Load-Time Decision}
\label{sec:hnsw-tuning}
%% ------------------------------------------------------------

The problem HNSW tuning solves is not recall alone --- it is the
\emph{relationship} between recall, memory consumption, index build
time, and query latency, all of which move when any single parameter
changes. The naive approach is to leave the defaults untouched
(\code{m=16}, \code{ef\_construct=100}) and compensate with a larger
\code{limit} at query time. This fails at scale because a sparser
graph with a wider beam wastes latency budget navigating dead-end
paths rather than improving candidate quality. The principle is that
index geometry must be specified explicitly at collection creation,
tuned to the cardinality and recall requirements of the workload, and
treated as a deployment-time configuration decision rather than a
default.

\begin{table}[h]
\centering
\caption{HNSW parameter guide for agent memory workloads}
\label{tab:hnsw-params}
\begin{tabularx}{\textwidth}{llX}
\toprule
\textbf{Parameter} & \textbf{Range} & \textbf{When to increase} \\
\midrule
\code{m} & 16--64 & Recall is below target at high collection sizes ($>$500K points); cosine text embeddings with high dimensionality. Cost: RAM grows linearly with \code{m}. \\
\code{ef\_construct} & 100--200 & Graph quality matters more than ingestion speed; this is a one-time cost paid per segment build. Cost: slower index construction. \\
\code{ef} (runtime) & 50--100 & Query recall is below target at serving time; latency budget allows a wider beam. Overridable per query via \code{SearchParams}. \\
\code{full\_scan\_threshold} & 1\,000--20\,000 & Reduce to force HNSW on smaller warm-up segments; increase to defer HNSW construction on tiny collections. \\
\bottomrule
\end{tabularx}
\end{table}

A production agent memory collection accumulates points continuously
as agents process conversations. This means index construction happens
incrementally, and the most expensive mistake is triggering repeated
partial segment rebuilds during a large initial load. Qdrant provides a
mechanism to suppress indexing during bulk ingestion: setting
\code{indexing\_threshold=0} in \code{OptimizersConfigDiff} disables HNSW
segment construction entirely until the threshold is restored. The correct
protocol is to create the collection, create payload indexes
(Section~\ref{sec:payload-indexes}), disable indexing, run the bulk upsert,
and then re-enable indexing so that one large HNSW segment is built from the
complete dataset rather than many small ones being merged.

\begin{lstlisting}[language=Python,
  caption={Collection creation with explicit HNSW config and bulk-load helpers},
  label={lst:collection-setup}]
# src/vector_store/collection_setup.py
from qdrant_client import QdrantClient, models

COLLECTION = "agent_memory"

def create_agent_memory_collection(client: QdrantClient, dims: int) -> None:
    client.create_collection(
        collection_name=COLLECTION,
        vectors_config=models.VectorParams(
            size=dims,
            distance=models.Distance.COSINE,
        ),
        hnsw_config=models.HnswConfigDiff(
            m=32,
            ef_construct=200,
            full_scan_threshold=5_000,
            max_indexing_threads=0,  # 0 = use all available cores
        ),
        optimizers_config=models.OptimizersConfigDiff(
            indexing_threshold=10_000,
        ),
    )

def disable_indexing_for_bulk_load(client: QdrantClient) -> None:
    client.update_collection(
        collection_name=COLLECTION,
        optimizers_config=models.OptimizersConfigDiff(indexing_threshold=0),
    )

def enable_indexing_after_load(client: QdrantClient) -> None:
    # MAINTENANCE WINDOW: triggers background HNSW segment construction.
    # Do not call during live traffic; route requests away from this node first.
    client.update_collection(
        collection_name=COLLECTION,
        optimizers_config=models.OptimizersConfigDiff(indexing_threshold=10_000),
    )
\end{lstlisting}

Two decisions in Listing~\ref{lst:collection-setup} deserve explanation.
\code{max\_indexing\_threads=0} instructs Qdrant to use all available CPU
cores during segment construction --- this is safe for a dedicated service but
should be capped to a lower value (e.g., \code{2}) if the Qdrant node also
serves live traffic. The \code{disable\_indexing\_for\_bulk\_load} and
\code{enable\_indexing\_after\_load} pair deliberately uses two separate
functions rather than a context manager because re-enabling indexing is a
maintenance-window operation: calling \code{update\_collection} with a new
\code{indexing\_threshold} triggers background HNSW segment construction that
will increase CPU and memory pressure for the duration of the build. The
separation makes the maintenance boundary explicit in calling code.

\begin{warningbox}
Calling \code{update\_collection} with a new \code{hnsw\_config} --- not
just \code{optimizers\_config} --- triggers a full index rebuild. The official
Qdrant documentation states: \emph{``We recommend against using this in a
production database as it may introduce huge overhead due to the rebuilding
of the index.''} Any change to \code{m} or \code{ef\_construct} on an
existing collection must be performed in a maintenance window with live
traffic routed away from the node.
\end{warningbox}

With index geometry established, the next irreversible decision must be made
before a single vector is written: the distance metric.


%% ------------------------------------------------------------
\section{Distance Metric Is Permanent}
\label{sec:metric-commitment}
%% ------------------------------------------------------------

The problem metric selection solves is matching the mathematical operation
of similarity comparison to the mathematical space in which the encoder
was trained. The naive approach --- defaulting to cosine for all embeddings ---
fails when the encoder was trained using dot product as its objective,
because the distances Qdrant computes will not correspond to the semantic
distances the encoder learned. The principle is that metric choice must
follow encoder training objective, and it must be declared once, permanently.

\begin{warningbox}
\textbf{The distance metric is locked at collection creation. There is no
in-place migration path.} Changing the metric requires deleting the collection,
re-creating it with the new metric, and re-ingesting every vector from source.
On a collection with tens of millions of agent memories, this is a multi-hour
operation. Choose the correct metric before writing the first point.
\end{warningbox}

Three metrics cover the production cases for text agent memory. \code{COSINE}
is appropriate for raw, unnormalised text embeddings such as those produced by
OpenAI's \code{text-embedding-3-small} and \code{text-embedding-3-large} or
\code{sentence-transformers} models --- the metric normalises vectors at query
time, making it invariant to embedding magnitude. \code{DOT} is appropriate
when the encoder was trained with dot product as its objective and guarantees
unit-norm output, such as Cohere \code{embed-v3} or \code{bge-*} models
configured with unit normalisation: because the vectors are already
unit-normalised, dot product and cosine are mathematically identical, but
\code{DOT} avoids the redundant normalisation step inside Qdrant's scoring
kernel, yielding a modest throughput improvement at high query volumes.
\code{EUCLID} is appropriate for encoders whose training objective was
Euclidean distance --- typically geometric embedding models applied to
structured coordinate data, not free-text memory content.

\begin{lstlisting}[language=Python,
  caption={Metric selection at collection creation for two common encoder families},
  label={lst:metric-selection}]
# src/vector_store/collection_setup.py (continued)
from qdrant_client import QdrantClient, models

# Raw OpenAI embeddings: magnitude varies; cosine is correct.
def create_openai_collection(client: QdrantClient) -> None:
    client.create_collection(
        collection_name="agent_memory_openai",
        vectors_config=models.VectorParams(
            size=1536,
            distance=models.Distance.COSINE,
        ),
    )

# Cohere embed-v3 with int8 output: unit-norm, pre-quantised.
# Requires qdrant-client >= 1.9.0 and Qdrant server >= 1.9.0.
def create_cohere_collection(client: QdrantClient) -> None:
    client.create_collection(
        collection_name="agent_memory_cohere",
        vectors_config=models.VectorParams(
            size=1024,
            distance=models.Distance.DOT,
            datatype=models.Datatype.UINT8,  # int8 compact storage
        ),
    )
\end{lstlisting}

Two decisions in Listing~\ref{lst:metric-selection} deserve explanation.
\code{models.Datatype.UINT8} is only valid when the encoder produces
pre-quantised 8-bit integer output --- using it with a float32 encoder will
silently truncate values and destroy recall. The two functions are deliberately
named after their encoder family, not their collection name, to make the
encoder-metric coupling explicit in calling code and prevent future engineers
from reusing the collection-creation helper with a different encoder.

\begin{notebox}
The choice between \code{COSINE} and \code{DOT} for a new encoder family is
not guesswork: read the encoder model card. If the card specifies ``cosine
similarity'' as the recommended scoring method, use \code{COSINE}. If it
specifies ``dot product'' or guarantees unit-norm output, use \code{DOT}. If
the card is absent or ambiguous, default to \code{COSINE} --- it is never
incorrect for text, only occasionally slightly slower than \code{DOT}.
\end{notebox}

With collection geometry and metric committed, the collection is structurally
ready but not yet queryable with filters --- that requires the payload indexes
that must be installed before any data arrives.


%% ------------------------------------------------------------
\section{Payload Indexes Must Precede Data}
\label{sec:payload-indexes}
%% ------------------------------------------------------------

The problem this section solves is making structured metadata on stored
points efficiently queryable without scanning the entire collection on every
filtered search. The naive approach --- running searches with payload filters
on unindexed fields --- does not raise an error; it silently falls back to a
full collection scan, consuming proportionally more CPU and memory the larger
the collection grows. By the time the performance problem is noticed in
production, the collection has hundreds of millions of points and rebuilding
the HNSW index to incorporate the new payload indexes is a multi-hour
maintenance operation. The principle is that payload indexes are created
immediately after collection creation and before any data ingestion, and the
setup function is written to be idempotent so it can be called safely on
every deployment.

The five fields that every agent memory point must carry --- and that must be
indexed --- are: \code{user\_id} (tenant identification), \code{memory\_type}
(episodic, semantic, or procedural), \code{importance\_score} (float for
range-based relevance filtering), \code{created\_at} (UNIX integer timestamp,
principal-indexed for time-range queries), and \code{ttl\_expires\_at} (UNIX
integer timestamp for TTL sweep range scans). Both timestamp fields are stored
as UNIX integers rather than ISO datetime strings. While Qdrant does support
\code{is\_principal} on \code{DatetimeIndexParams} as of v1.11.0, storing
timestamps as integers with \code{IntegerIndexParams} offers an additional
advantage: the parameterized integer index exposes separate \code{lookup} and
\code{range} flags, allowing lookup indexes to be disabled on fields that are
queried exclusively by range comparators. This reduces the memory footprint of
the integer index without any functional cost.

\begin{notebox}
Qdrant's official documentation states: \emph{``Payload indexes should be
created before ingesting data. Qdrant's filterable HNSW index only benefits
from additional filter-aware edges when it is generated after the payload
indexes have been created. If you create a payload index after data has
already been ingested, you need to rebuild the HNSW index to take advantage
of the new payload indexes.''} This is not a recommendation --- it is a
structural requirement.
\end{notebox}

\begin{lstlisting}[language=Python,
  caption={Idempotent payload index setup loop},
  label={lst:payload-indexes}]
# src/vector_store/index_setup.py
import logging
from qdrant_client import QdrantClient, models
from qdrant_client.http.exceptions import UnexpectedResponse

log = logging.getLogger(__name__)

PAYLOAD_INDEX_SPECS = {
    "user_id": models.KeywordIndexParams(
        type=models.KeywordIndexType.KEYWORD,
        is_tenant=True,
    ),
    "memory_type": models.KeywordIndexParams(
        type=models.KeywordIndexType.KEYWORD,
        is_tenant=False,
    ),
    "importance_score": models.FloatIndexParams(
        type=models.FloatIndexType.FLOAT,
        is_principal=False,  # not a primary ordering axis
    ),
    "created_at": models.IntegerIndexParams(
        type=models.IntegerIndexType.INTEGER,
        is_principal=True,   # primary time-range ordering axis
        lookup=False,        # no equality queries on timestamps
        range=True,
    ),
    "ttl_expires_at": models.IntegerIndexParams(
        type=models.IntegerIndexType.INTEGER,
        is_principal=False,
        lookup=False,
        range=True,
    ),
}

def setup_payload_indexes(client: QdrantClient, collection: str) -> None:
    for field, schema in PAYLOAD_INDEX_SPECS.items():
        try:
            client.create_payload_index(
                collection_name=collection,
                field_name=field,
                field_schema=schema,
                wait=True,
            )
        except UnexpectedResponse as exc:
            if exc.status_code == 400:
                log.info("Index already exists for %r, skipping.", field)
            else:
                raise
\end{lstlisting}

Three decisions in Listing~\ref{lst:payload-indexes} deserve explanation.
The \code{try/except UnexpectedResponse} block with an HTTP 400 guard makes
the function idempotent: Qdrant returns a 400 status when an index already
exists with the same schema, so this function can run on every deployment
without failing on collections that were already configured. Errors with any
other status code are re-raised immediately to surface genuine problems.
The \code{lookup=False, range=True} combination on both timestamp fields is
intentional: these fields are never queried with equality match conditions ---
only with range comparators. Disabling \code{lookup} reduces the memory
footprint of the integer index without any functional cost. Finally,
\code{wait=True} on every call ensures the function blocks until each index
is confirmed built before returning; calling code can then proceed to
ingestion without a race condition.

\begin{tipbox}
After deploying to production, enable strict mode on the collection by
setting \code{unindexed\_filtering\_retrieve=false}. With strict mode
active, Qdrant returns an error if a search query filters on a field
without a payload index, turning silent performance degradation into a
visible configuration error at the API boundary.

Full-text indexing on the \code{content} field via \code{TextIndexParams}
and \code{MatchText} conditions is deferred to Chapter~12 \S12.3 hybrid
retrieval.
\end{tipbox}

With all payload indexes registered, the collection is ready for data
ingestion --- and the reason those indexes had to come first will become
clear in the next section, which explains how Qdrant's query planner uses
them during HNSW traversal.


%% ------------------------------------------------------------
\section{Recall Does Not Collapse When Filters Are Index-Aware}
\label{sec:filter-aware-hnsw}
%% ------------------------------------------------------------

The problem this section solves is the recall collapse that occurs when
filtered vector search is implemented as two independent steps: run the
approximate nearest-neighbour search first, then apply a payload predicate
to the results. The naive approach treats the ANN search and the filter
as orthogonal operations, which works correctly only when the filter
passes the majority of points in the collection. When the filter is
selective --- as it always is in a multi-tenant collection where
\code{user\_id} identifies one user among tens of thousands --- the top-K
candidates returned by ANN are almost entirely from other tenants; the
post-filter step reduces the result set to near-zero without increasing
recall. The principle is that the filter must participate in the graph
traversal, not follow it.

Qdrant achieves this by using payload indexes as cardinality oracles during
query planning. When a search query arrives with a \code{query\_filter},
Qdrant's planner reads the estimated number of matching points from the
relevant payload index and uses it to choose a search strategy. When the
filter is loose (many matching points), the planner runs
\emph{filtered HNSW traversal}: the beam search navigates the graph and
evaluates the filter condition at each hop using the payload index for fast
lookup. When the filter is highly selective (very few matching points), the
planner runs a \emph{payload scan followed by ANN reranking}: it reads only
the indexed matching points and reranks them by vector similarity.

These two strategies have different relationships to the build-time ordering
rule from Section~\ref{sec:payload-indexes}. The filtered HNSW traversal path
benefits from payload indexes existing at segment build time, because
filter-aware edges are woven into the graph during initial construction.
The payload scan path uses the index at query time and does not depend on
build-time ordering. In both cases, however, the index must exist --- and in
both cases the post-filter antipattern is avoided.

\begin{lstlisting}[language=Python,
  caption={Anti-pattern and correct filter-aware search side by side},
  label={lst:filter-search}]
# src/vector_store/memory_search.py
from qdrant_client import QdrantClient, models

# ANTI-PATTERN: ANN first, Python-side filter after.
# Recall collapses when user_id matches < 1% of collection.
def search_naive(client, collection, vector, user_id, limit):
    all_hits = client.search(collection, query_vector=vector, limit=500)
    return [h for h in all_hits if h.payload["user_id"] == user_id][:limit]

# CORRECT: native query_filter; planner chooses traversal strategy.
def search_memory(
    client: QdrantClient,
    collection: str,
    vector: list[float],
    user_id: str,
    min_importance: float,
    limit: int = 10,
    hnsw_ef: int = 80,
) -> list[models.ScoredPoint]:
    filt = models.Filter(
        must=[
            models.FieldCondition(
                key="user_id",
                match=models.MatchValue(value=user_id),
            ),
            models.FieldCondition(
                key="importance_score",
                range=models.Range(gte=min_importance),
            ),
        ]
    )
    return client.search(
        collection_name=collection,
        query_vector=vector,
        query_filter=filt,          # <- search() uses query_filter
        limit=limit,
        with_payload=True,
        with_vectors=False,
        search_params=models.SearchParams(hnsw_ef=hnsw_ef),
    )
\end{lstlisting}

Three decisions in Listing~\ref{lst:filter-search} deserve explanation.
The anti-pattern function requests 500 results to compensate for the
expected post-filter discard rate --- this is the pattern to eliminate, not
to refine. The \code{with\_vectors=False} flag halves the wire payload for
every response: agent memory retrieval uses the scored content, not the raw
embedding. The \code{models.SearchParams(hnsw\_ef=hnsw\_ef)} override applies
a per-query beam width, enabling the calling code to tune recall at query time
without touching the collection-level \code{ef} configuration.

\begin{warningbox}
The \code{search()} method accepts the parameter \code{query\_filter}.
The \code{scroll()} method accepts \code{scroll\_filter}. These are
different keyword argument names on different methods. Passing
\code{query\_filter=} to a \code{scroll()} call may not raise a
\code{TypeError} if the extra keyword is silently absorbed; the filter
will be ignored and the scroll will return unfiltered results. Always
use the correct named parameter explicitly for each method.
\end{warningbox}

With filter-aware traversal understood, the design is ready to be applied to
the multi-tenant isolation problem --- where the \code{user\_id} filter is
not an occasional constraint but the invariant present in every query.


%% ------------------------------------------------------------
\section{One Collection Isolates All Tenants}
\label{sec:multi-tenancy}
%% ------------------------------------------------------------

The problem this section solves is maintaining strict per-user memory
isolation in a system with a large and growing number of users, without
multiplying infrastructure overhead with the user count. The naive
approach --- creating one Qdrant collection per user --- appears to provide
the strongest isolation because the data boundaries are physical. It fails
because each Qdrant collection carries fixed overhead: a write-ahead log,
optimizer threads, HNSW segment files, and memory-mapped payload storage.
The official Qdrant documentation states directly: \emph{``In most cases,
you should only use a single collection with payload-based partitioning.''}
At ten thousand users, ten thousand collections exhaust file descriptors
and WAL buffer allocations on any realistic server; performance degrades
non-linearly before that limit is reached. The principle is that isolation
is a query-time property enforced by mandatory payload filters, not a
storage-time property enforced by collection boundaries.

The single-collection approach works because of the \code{is\_tenant=True}
flag introduced in Section~\ref{sec:payload-indexes}. This flag instructs
Qdrant's storage engine to physically co-locate data belonging to each
tenant on disk, reducing disk reads required for a tenant-scoped query.
Combined with filter-aware HNSW traversal, the recall and latency
characteristics of a per-user query against the shared collection are
equivalent to querying a per-user collection --- without the operational cost.

Isolation is logical, not physical. The application layer is responsible
for ensuring that every query includes the correct \code{user\_id} filter.
Chapter~16 covers the FastAPI middleware that extracts the \code{user\_id}
claim from the JWT and passes it into the query path; this section
establishes the data model and the filter factory that the middleware will
call.

\begin{lstlisting}[language=Python,
  caption={MemoryPayload model and tenant filter factory},
  label={lst:tenant-models}]
# src/vector_store/models.py
from pydantic import BaseModel
from typing import Literal

class MemoryPayload(BaseModel):
    user_id: str                                             # keyword, is_tenant=True
    memory_type: Literal["episodic", "semantic", "procedural"]  # keyword
    content: str
    importance_score: float                                  # float index
    created_at: int                                          # UNIX timestamp, principal
    ttl_expires_at: int | None                               # UNIX timestamp, range
\end{lstlisting}

\begin{lstlisting}[language=Python,
  caption={Tenant-scoped filter factory and search helper},
  label={lst:tenant-ops}]
# src/vector_store/tenant_ops.py
from qdrant_client import QdrantClient, models
from src.vector_store.models import MemoryPayload

def build_user_filter(
    user_id: str,
    memory_types: list[str] | None = None,
    min_importance: float | None = None,
) -> models.Filter:
    conditions: list[models.Condition] = [
        models.FieldCondition(
            key="user_id",
            match=models.MatchValue(value=user_id),
        )
    ]
    if memory_types:
        conditions.append(
            models.FieldCondition(
                key="memory_type",
                match=models.MatchAny(any=memory_types),
            )
        )
    if min_importance is not None:
        conditions.append(
            models.FieldCondition(
                key="importance_score",
                range=models.Range(gte=min_importance),
            )
        )
    return models.Filter(must=conditions)

def search_user_memories(
    client: QdrantClient,
    collection: str,
    query_vector: list[float],
    user_id: str,
    limit: int = 10,
    memory_types: list[str] | None = None,
    min_importance: float | None = None,
) -> list[models.ScoredPoint]:
    return client.search(
        collection_name=collection,
        query_vector=query_vector,
        query_filter=build_user_filter(user_id, memory_types, min_importance),
        limit=limit,
        with_payload=True,
        with_vectors=False,
    )
\end{lstlisting}

Two decisions in Listings~\ref{lst:tenant-models} and~\ref{lst:tenant-ops}
deserve explanation. The \code{conditions} list is annotated as
\code{list[models.Condition]} --- the correct supertype --- rather than
\code{list[models.FieldCondition]}: this allows callers to append
\code{models.IsNullCondition} or other condition subtypes without a type
error. \code{models.MatchAny(any=memory\_types)} is used for
\code{memory\_type} rather than multiple \code{MatchValue} conditions
wrapped in a \code{should} clause: \code{MatchAny} is a single condition
object that the planner evaluates in one index lookup, whereas a \code{should}
clause with multiple conditions requires the planner to union multiple index
reads and deduplicate the result set.

\begin{notebox}
The \code{MemoryPayload} model stores timestamps as \code{int} (UNIX epoch
seconds). This choice is intentional: the parameterized \code{IntegerIndexParams}
exposes separate \code{lookup} and \code{range} flags that allow disabling
the lookup component on pure range-query fields, reducing index memory
footprint. The \code{DatetimeIndexParams} type supports \code{is\_principal}
as of Qdrant v1.11.0 but does not expose the \code{lookup}/\code{range}
parameterization, making the integer representation preferable for this
workload.
\end{notebox}

The collection, indexes, and isolation pattern are now complete. The final
section connects them to the mem0 OSS memory layer introduced in Chapter~8.


%% ------------------------------------------------------------
\section{mem0 Connects to a Pre-Configured Collection}
\label{sec:mem0-backend}
%% ------------------------------------------------------------

The problem this section solves is wiring the mem0 OSS \code{Memory} class
--- which provides the agent-facing memory API used throughout this book ---
to the Qdrant collection configured in Sections~\ref{sec:hnsw-tuning}
through \ref{sec:multi-tenancy}, without allowing mem0 to silently overwrite
that configuration. The naive approach is to call \code{Memory.from\_config()}
with a Qdrant backend config and assume mem0 will use the existing collection.
This fails because mem0 creates the collection if it does not exist, using
default HNSW parameters --- and if the collection does not yet exist at the
time \code{from\_config()} is called (which is likely during first deploy),
it will be created with \code{m=16}, \code{ef\_construct=100}, and no payload
indexes. The principle is that collection existence and configuration are
asserted before the \code{Memory} object is constructed.

\begin{notebox}
\code{Memory} (OSS, imported as \code{from mem0 import Memory}) and
\code{MemoryClient} (Platform API, imported as
\code{from mem0 import MemoryClient}) are different classes.
\code{MemoryClient} connects to the hosted mem0 API and does not accept
a Qdrant backend config. Every reference to mem0 in this chapter uses
\code{Memory} --- the self-hosted OSS class that accepts the
\code{vector\_store} config dict shown below.
\end{notebox}

Chapter~8 Section~8.2 introduced \code{Memory.from\_config()} as the
entry point for the mem0 memory layer and showed the structure of the
config dict. Here, the \code{vector\_store} block in that config must carry
the collection name and embedding dimensions that exactly match what was
used in \code{create\_agent\_memory\_collection()}. A mismatch in
\code{embedding\_model\_dims} between the config and the
\code{VectorParams.size} set at collection creation will cause ingestion
failures or silent dimension errors at query time. The
\code{QdrantSettings} Pydantic model is the single source of truth for
both values.

\begin{lstlisting}[language=Python,
  caption={QdrantSettings and collection guard},
  label={lst:qdrant-settings}]
# src/memory/mem0_client.py
from pydantic import AnyHttpUrl, SecretStr
from pydantic_settings import BaseSettings

class QdrantSettings(BaseSettings):
    collection_name: str = "agent_memory"
    url: AnyHttpUrl
    api_key: SecretStr | None = None  # set QDRANT_API_KEY for Qdrant Cloud
    embedding_dims: int = 1536        # must match VectorParams.size

    model_config = {"env_prefix": "QDRANT_"}
\end{lstlisting}

\begin{lstlisting}[language=Python,
  caption={Collection guard, mem0 config builder, and Memory factory},
  label={lst:mem0-client}]
# src/memory/mem0_client.py (continued)
from qdrant_client import QdrantClient
from mem0 import Memory  # OSS class, NOT MemoryClient (Platform)
from src.vector_store.collection_setup import create_agent_memory_collection
from src.vector_store.index_setup import setup_payload_indexes

def ensure_collection_ready(
    client: QdrantClient,
    settings: QdrantSettings,
) -> None:
    if not client.collection_exists(settings.collection_name):
        create_agent_memory_collection(client, settings.embedding_dims)
        setup_payload_indexes(client, settings.collection_name)

def build_mem0_config(settings: QdrantSettings, base_config: dict) -> dict:
    # base_config carries "embedder" and "llm" blocks from Ch.8 §8.2
    return {
        **base_config,
        "vector_store": {
            "provider": "qdrant",
            "config": {
                "collection_name": settings.collection_name,
                "embedding_model_dims": settings.embedding_dims,
                # Cross-check key name against installed mem0ai version;
                # it has changed across minor releases.
                "url": str(settings.url),
                "api_key": settings.api_key.get_secret_value()
                           if settings.api_key else None,
            },
        },
    }

def get_memory(
    client: QdrantClient,
    settings: QdrantSettings,
    base_config: dict,
) -> Memory:
    ensure_collection_ready(client, settings)
    return Memory.from_config(build_mem0_config(settings, base_config))
\end{lstlisting}

Three decisions in Listings~\ref{lst:qdrant-settings}
and~\ref{lst:mem0-client} deserve explanation. \code{ensure\_collection\_ready}
checks existence with \code{client.collection\_exists()} before calling the
creation and index setup functions --- this is the guard that prevents mem0
from racing with application startup and creating an unconfigured collection.
\code{settings.api\_key.get\_secret\_value()} unwraps the \code{SecretStr}
explicitly rather than relying on implicit string conversion: Pydantic's
\code{SecretStr} does not expose its value via \code{str()}, which would
produce \code{"**********"} and cause a silent authentication failure against
Qdrant Cloud. The \code{build\_mem0\_config} function accepts a
\code{base\_config} dict carrying the \code{embedder} and \code{llm} blocks
introduced in Chapter~8 \S8.2, merging them with the \code{vector\_store}
block via \code{**base\_config} so that all three required sections are
present before \code{Memory.from\_config()} is called.

\begin{tipbox}
Cross-check the \code{embedding\_model\_dims} key name against the version
of \code{mem0ai} installed in your environment. The key has changed between
minor releases. Confirm the correct name by inspecting \code{mem0.memory.main}
or reading the changelog for the pinned release. Pin \code{mem0ai} in
\code{pyproject.toml} to avoid silent key-name drift across environments.
\end{tipbox}

Chapter~11 builds directly on the collection established here, adding scalar
quantization to compress the vector storage footprint --- an optimisation that
interacts with the HNSW \code{m} parameter tuned in
Section~\ref{sec:hnsw-tuning} and that requires the same
\code{ensure\_collection\_ready} guard introduced in this section.
